\section{Related Work}


\textbf{Compiler IR for Approximations.}
EnerJ~\cite{enerJ} supports approximate computing via type declarations that mark data as precise or approximate, enabling static analysis to isolate precise and approximate code; hardware-level approximations include reduced DRAM refresh rates, lower SRAM voltage, and approximate floating-point operations. While effective for general-purpose applications, EnerJ’s design limits its applicability to domain-specific workloads such as HDC and does not support heterogeneous accelerator targets. Rely~\cite{rely} provides a programming language and analysis framework that enables developers to specify quantitative reliability requirements for programs running on unreliable hardware (i.e., hardware where arithmetic operations or memory accesses may fail with some probability). ApproxHPVM~\cite{approxhpvm} and ApproxTuner~\cite{approxtuner} provide compiler- and runtime-level support for accuracy-aware approximations in deep learning and image processing. ApproxHPVM encodes per-operation error metrics in the HPVM IR and maps computations to approximate hardware to meet end-to-end quality constraints. %, achieving up to 9× speedup and 11× energy reduction. 
ApproxTuner extends this with a three-stage tuning framework that explores approximation-performance trade-offs at development, install, and runtime, achieving significant speedups with minimal accuracy loss. Both frameworks are effective for functional, tensor-based workloads but do not directly support domain-specific paradigms like hyperdimensional computing (HDC), which require more intrusive approximation techniques to maintain acceptable quality-of-service given their inexpensive, non-backpropagation-based learning algorithms. %Moreover, HDC applications rely on specialized HDC operations interleaved with general-purpose parallel C++ code, making domain-specific compiler and approximation support essential.


%ApproxHPVM
%Rely
%TVM
%enerj

\textbf{Hyperdimensional Computing Approximation.}
%MicroHD~\cite{microhd} is an accuracy-driven optimization framework for hyperdimensional computing targeting TinyML and edge environments. It explores a limited hyperparameter-based approximation space, including encoding dimensionality, quantization bit-width, number of hypervectors, and encoding algorithm choice, which is smaller compared to frameworks that support more diverse or fine-grained approximations. MicroHD employs a combination of binary search and greedy algorithms to efficiently select configurations that minimize memory and compute requirements while satisfying user-specified accuracy constraints, achieving up to 200–$266\times$ resource savings with less than 1\% accuracy loss. The framework is implemented using TorchHD~\cite{torchhd}, a Python library for HDC, and does not leverage a compiler IR; as a result, all transformations and approximations must be manually implemented, and it lacks flexibility in defining application-specific tuning parameters or hardware-specific knobs. 
MicroHD~\cite{microhd} is an accuracy-driven HDC optimization framework for TinyML and edge settings that explores a limited hyperparameter-based approximation space (e.g., dimensionality, quantization, and encoding choices). It uses binary search and greedy methods to meet accuracy constraints while reducing resource usage, achieving up to 200–266× savings with <1\% accuracy loss. Built on TorchHD~\cite{torchhd} without a compiler IR, it requires manual implementation of approximations and offers limited flexibility for application- or hardware-specific tuning.
In contrast, our approach is extensible, as it automatically constructs the approximation space, allows the definition of custom application-specific tuning parameters with ease, and supports hardware approximation knobs not previously available, while also providing compiler-based insertion and management of software- and hardware-level approximations across heterogeneous targets.~\cite{kazemi2022achieving} show that HDC can tolerate hardware-level approximation using multi-bit FeFET in-memory computing, achieving software-equivalent accuracy with up to $826\times$ energy and 30× latency improvements. This demonstrates that HDC is resilient to approximate hardware execution, motivating compiler-based frameworks to systematically exploit such hardware knobs for performance and efficiency gains.


%% No need for the following. we are not proposing an autotuner.
%\textbf{Approximation Tuning Frameworks.} OpenTuner~\cite{opentuner} is an extensible autotuning framework that automates exploration of large configuration spaces using an ensemble of search techniques. It allows developers to define custom search spaces, objectives, and search algorithms, including a multi-armed bandit strategy to allocate more evaluations to promising search techniques while reducing time spent on underperforming ones. This approach enables efficient optimization across diverse domains and hardware targets. OpenTuner has been applied to performance tuning, multi-objective optimization, and compiler parameter selection, providing a flexible foundation for frameworks that require systematic exploration of program parameters, including in approximate or domain-specific computing. pyATF~\cite{pyatf} is a Python-based auto-tuning framework that supports constraint-aware parameter spaces, enabling efficient pruning of invalid configurations. It employs multi-armed bandit algorithms to allocate search effort to promising strategies, similar to OpenTuner, but with built-in handling of interdependent parameters. Compared to OpenTuner, pyATF is better suited for domains with complex parameter relationships, such as compiler flags or hardware knobs, while OpenTuner offers more flexible, general-purpose search and custom search heuristics.






